\documentclass[10pt,a4paper]{amsart}
\usepackage[margin=19mm]{geometry}
\usepackage{amsmath,amssymb,mathtools}
\usepackage[scr=rsfs,cal=euler]{mathalfa}
\usepackage{xfrac}
\usepackage[unicode,hidelinks,bookmarks=false]{hyperref}
\newcommand{\ie}{i.e.\ }
\newcommand{\cf}{cf.\ }
\newcommand{\resp}{respectively\ }
\newcommand{\Subsection}{\subsection}
\providecommand{\nobreakdash}{\nobreak\hbox{-}}
\setlength{\emergencystretch}{2em}
\multlinegap0pt
\usepackage{amssymb}
\usepackage{mathtools}
\newtheorem{lemma}{Lemma}
\newtheorem{theorem}{Theorem}
\newtheorem{proposition}{Proposition}
\newcommand{\F}{\mathbb F}
\newcommand{\GG}{G}
\newcommand{\Z}{\mathbb Z}
\newcommand{\dd}{\mathrm d}
\newcommand{\ee}{\mathrm e}
\newcommand{\pp}{\mathfrak p}
\title{Integer group determinants for the elementary abelian group of order 49}
\author{Chatchawan Panraksa}
\date{Source: arXiv:2609.09542v1 (CC BY 4.0)}
\begin{document}
\maketitle

\noindent\emph{Source and licence.} Integer group determinants for the elementary abelian group of order 49. Version arXiv:2609.09542v1 (CC BY 4.0), distributed by arXiv under the Creative Commons Attribution 4.0 International licence, http://creativecommons.org/licenses/by/4.0/, as recorded in the arXiv licence metadata for this exact version and shown on the abstract page https://arxiv.org/abs/2609.09542v1. Full licence notice: \texttt{licenses/arxiv-2609.09542-CC-BY-4.0.txt}.\par\smallskip

\noindent\textit{Editorial scope.} Counted unit: the article's theorem thm:crit11 (source0.tex lines 940--959). The article's complete original proof of it is delivered with it (source0.tex lines 961--1034, one \texttt{proof} environment of 74 lines). No mathematics is written, repaired or re-derived: the statement and the proof are the article's own lines, in the article's own notation, and no retained line is substituted or re-flowed.  Retained uncounted as the source-local context the proof needs: lines 404--417; lines 475--483; lines 743--782; lines 890--900; lines 1189--1193.  Nothing was omitted from any retained span, so the package carries no omission record.  No retained line contains a citation, so the wrapper carries no bibliography and no bibliography entry.  No retained line contains a figure environment, an image-inclusion command or drawing code, so no image is delivered and the package declares no asset dependency.  Editorial formatting only, in addition: an AMS \texttt{amsart} preamble that restates the article's own declarations and macros used by the retained lines, namely the environments \texttt{lemma}, \texttt{theorem}, \texttt{proposition} on separate counters, together with the standard packages those lines need.  The retained environments print in this document as Lemma 1, Theorem 1, Proposition 1 in order of appearance, whereas the article prints the same items with the numbering of its own full document.  The complete native source of the article as distributed in the arXiv e-print for this exact version is bundled unchanged as \texttt{sources/209826/source0.tex}.
\begin{lemma}[Reconstruction]\label{lem:interp}
Let $c^{(r)}\in\Z^7$ have common coordinate sum $s$, and put
\[
 F(i,j)=c^{(\infty)}(j)+\sum_{b\in\F_7}c^{(b)}(i+bj).
\]
\begin{enumerate}
\item[(i)] If the $c^{(r)}$ are the ray vectors of $v$, then
$F(i,j)=7v_{i,j}+s$.
\item[(ii)] Given $f_r\in R$ with $f_r\equiv s\pmod\pp$, take their
ray vectors of sum $s$. An integral vector with character data
$(s;\sigma_t(f_r))$ exists if and only if $F\equiv s\pmod7$ on
$\F_7^2$. When it exists, it is $v=(F-s)/7$.
\end{enumerate}
\end{lemma}

\begin{lemma}\label{lem:identity}
Let $x_3,x_4\in\F_7[b]$ have degrees at most $3,2$, respectively,
and write $q_i=[b^i]x_4$ and $x_{33}=[b^3]x_3$. Then
\begin{align*}
 \sum_b(-x_4(b)^3-2x_3(b)x_4(b))&=q_2^3,\\
 \sum_b b(-x_4(b)^3-2x_3(b)x_4(b))
     &=3q_1q_2^2+2x_{33}q_2.
\end{align*}
\end{lemma}

\begin{theorem}[Ray residues of unit multiples]\label{thm:twisted}
Fix $y\in R$ coprime to $\pi$, and put $d=\dd(y)$, $e=\ee(y)$.
As $u$ ranges over $R^\times$, the ray residues $\bar c_r(yu)$,
normalized to sum seven, have the following exact descriptions.
\begin{enumerate}
\item[(a)] On the slice $\lambda=1$, the residues $\bar c_1(yu)$
are precisely the solutions with $\gamma_6=0$, $\gamma_5=-1$ of
\begin{align*}
 \gamma_2&=-\gamma_4^3-2\gamma_3\gamma_4+d,\\
 \gamma_0&=1+3\gamma_4^5+5\gamma_3\gamma_4^3+4\gamma_1\gamma_4
                  +3d\gamma_4^2+2d\gamma_3+e.
\end{align*}
The coordinates $(\gamma_1,\gamma_3,\gamma_4)$ range freely over $\F_7^3$.
\item[(b)] The residues $\bar c_2(yu)$ are precisely the solutions with
$\gamma_6=\gamma_5=0$ and $\gamma_4\ne0$ of
\[
 \gamma_1=-\gamma_3^3\gamma_4^4
                -3\gamma_2\gamma_3\gamma_4^5+d\gamma_4.
\]
The coordinates $(\gamma_0,\gamma_2,\gamma_3,\gamma_4)$ range freely
over $\F_7^3\times\F_7^\times$.
\item[(c)] The residues $\bar c_3(yu)$ are precisely the solutions with
$\gamma_6=\gamma_5=\gamma_4=0$ and $\gamma_3\ne0$ of
\begin{equation}\label{eq:twisted-third}
 4\gamma_3^2(\gamma_0-1)+\gamma_1\gamma_2\gamma_3
              -\gamma_2^3-d\gamma_3^3=0.
\end{equation}
Equivalently,
\begin{equation}\label{eq:twisted-third-solved}
 \gamma_0=1+2d\gamma_3+2\gamma_2^3\gamma_3^{-2}
                              -2\gamma_1\gamma_2\gamma_3^{-1}.
\end{equation}
The coordinates $(\gamma_1,\gamma_2,\gamma_3)$ range freely over
$\F_7^2\times\F_7^\times$.  At every fixed
$(\gamma_2,\gamma_3)$ there are exactly seven residues.
\end{enumerate}
Thus the residue possibilities of $\pi^ryu$, for $r=1,2,3$, depend on
the ideal $(y)$ only through its defect pair; for $r=2,3$ only $d$ is
needed.
\end{theorem}

\begin{theorem}[The profile $(1^7,2)$]\label{thm:crit10}
For prescribed cofactor ideals and each $\kappa\in\F_7$, the profile
$(1^7,2)$ with generic leading residue $1$ and coordinate-sum class
$\kappa$ is realizable if and only if
\[
 D_0=\sum_b\dd(\mathfrak a_b)=\pm1\quad\text{in }\F_7.
\]
When this holds, every sum $s\equiv7(1-\kappa)\pmod{49}$ is
realizable. There is no condition on the second defects or on
the special cofactor ideal.
\end{theorem}

\begin{proposition}\label{prop:witness12}
For every $n\in\Z$ there is a vector $v\in\Z^{49}$ with
$s(v)=49n$ and orbit norms $(343,7,7,7,7,7,7,7)$. Consequently,
$7^{12}\Z\subseteq S(\GG)$.
\end{proposition}

\begin{theorem}[The second exceptional profile]\label{thm:crit11}
Fix the placed cofactor ideals, with generic leading residue $1$.
\begin{enumerate}
\item[(a)] The profile $(1^7,2)$ with coordinate sum of valuation
exactly two is realizable if and only if $D_0=\pm1$.
\item[(b)] The profile $(1^7,3)$ with coordinate-sum class
$\kappa\in\F_7$ is realizable if and only if $D_0=D_1=0$ and
\[
 \begin{cases}
 D_2\ne0, &\text{or}\\
 D_2=0,\ D_3+d_\infty\ne0,\ \kappa\ne1+E_0, &\text{or}\\
 D_2=0,\ D_3+d_\infty=0,\ \kappa=1+E_0.
 \end{cases}
\]
When feasible, every $s\equiv7(1-\kappa)\pmod{49}$ is realizable.
This profile has determinant valuation eleven for $\kappa\ne1$;
the case $\kappa=1$ gives the construction of $7^{12}\Z$ in
Proposition~\ref{prop:witness12}.
\end{enumerate}
\end{theorem}

\begin{proof}
For (a), apply Theorem~\ref{thm:crit10} with $\kappa=1$ and
$s=49u$, $7\nmid u$.

For (b), we first derive the necessary constant condition without
restricting the free coordinates. The special valuation three
forces the quadratic coefficient of $x_4$ to vanish.
Write $x_4(b)=q_0+qb$, and let
$x\ne0$ and $c$ be the cubic and quadratic coefficients of $x_3$.
The two identities in Lemma~\ref{lem:identity} give $D_0=D_1=0$.
Write $D=D_2$ and let $B=\sum_b b x_1(b)$ be the first moment
of $x_1$. The remaining nonconstant conditions give
\[
 \gamma_{\infty,3}=x,\qquad
 r:=\gamma_{\infty,2}=-2xq-D,\qquad
 \gamma_{\infty,1}=-B,\qquad \sum_bx_1(b)=0.
\]
The generic constant law of Theorem~\ref{thm:twisted}(a) sums to
\begin{equation}\label{eq:generic-constant}
 \sum_b\gamma_0(b)
 =-5xq^3+4qB+3q^2D+2cD+2xD_3+E_0.
\end{equation}
Indeed, the contributing sums are
$\sum x_3x_4^3=-xq^3$, $\sum x_1x_4=qB$,
$\sum d_bx_4^2=q^2D$ and $\sum d_bx_3=cD+xD_3$;
$\sum x_4^5=0$. These follow from $D_0=D_1=0$ and the usual
power sums in $\F_7$.
The special law in Theorem~\ref{thm:twisted}(c) gives
\[
 \gamma_{\infty,0}
 =1+2d_\infty x+2r^3x^{-2}+2Brx^{-1}.
\]
Adding the two constant laws and substituting $r=-2xq-D$, we
rewrite the constant condition as
\begin{equation}\label{eq:uniform-eleven}
 \kappa-1-E_0
 =2x(D_3+d_\infty)
  +2D\left(c+\frac{qD}{x}-\frac{D^2}{x^2}-\frac{B}{x}\right).
\end{equation}
Here the only denominator is a power of $x\ne0$. Put
$C=D_3+d_\infty$ and $K=\kappa-1-E_0$, the remaining prescribed
constant. If $D=0$, then the constant condition
\eqref{eq:uniform-eleven} reduces to $K=2xC$, so $K$ and $C$
must vanish together. This gives the necessary alternatives.

For sufficiency, we can set $q=B=0$ and solve for the remaining
coefficients $x,c$. Choose the generic free coordinates as
\[
 x_4(b)=x_1(b)=0,\qquad x_3(b)=xb^3+cb^2,
 \qquad x\ne0.
\]
Then $\gamma_2(b)=d_b$ and
$\gamma_0(b)=1+2d_b(xb^3+cb^2)+e_b$. Set
\[
 \gamma_{\infty,3}=x,\quad \gamma_{\infty,2}=-D,\quad
 \gamma_{\infty,1}=0,\quad
 \gamma_{\infty,0}=1+2d_\infty x-2D^3x^{-2}.
\]
These special coordinates satisfy Theorem~\ref{thm:twisted}(c).
Conditions (K$_1$)--(K$_4$) follow from
$D_0=D_1=0$ and the chosen polynomial degrees; (K$_5$) and
(K$_6$) hold as in the preceding proof. The constant condition is
\begin{equation}\label{eq:explicit-eleven}
 K=2xC+2cD-2D^3x^{-2}.
\end{equation}
If $D\ne0$, then setting $x=1$ determines
$c=D^2+(K-2C)/(2D)$.
If $D=0$ and $C\ne0$, then take $x=K/(2C)$ and $c=0$;
feasibility ensures $x\ne0$. If $D=C=0$, then take $x=1,c=0$,
since $K=0$.
Fullness in Theorem~\ref{thm:twisted} supplies the unit multiples,
and Lemma~\ref{lem:interp} reconstructs a vector for every sum
in the stated class.
\end{proof}

\end{document}
